\unit{Exposição 149 — Teoremas de dualidade para os feixes algébricos coerentes}
\sid{149}
\src{169}
\seminar{(Maio de 1957)}
\paperhead{TEOREMAS DE DUALIDADE PARA OS FEIXES ALGÉBRICOS COHERENTES}%
  {por Alexandre GROTHENDIECK}

Os resultados que seguem, inspirados no \fquote{teorema de dualidade algébrica}
de Serre, foram obtidos nos invernos de 1955 e 1956. Demonstram-se muito
simplesmente com o auxílio de resultados bastante elementares sobre a cohomologia dos espaços
projetivos [3] e mediante o uso intensivo da álgebra homológica de
Cartan--Eilenberg, na forma de [2].

\fsection{1}{Os Ext de feixes de módulos ({[2]}, caps. 3 e 4)}

Seja $X$ um espaço topológico provido de um feixe $\Osh$ de anéis com unidade
(não necessariamente comutativos). Consideramos a categoria abeliana
$\Ccat^{\Osh}$ dos feixes de $\Osh$-módulos, chamados também $\Osh$-Módulos.
Sabe-se que todo objeto desta categoria admite uma resolução injetiva, o que
permite definir nela os funtores $\operatorname{Ext}$ com as propriedades
formais bem conhecidas. Mais precisamente, para evitar confusões,
denotamos por $\operatorname{Hom}_{\Osh}(X;A,B)$, ou simplesmente por
$\operatorname{Hom}(X;A,B)$, o grupo abeliano dos $\Osh$-homomorfismos de $A$
em $B$, enquanto que $\Homsh_{\Osh}(A,B)$ denotará o feixe de germes
de homomorfismos de $A$ em $B$ ($A,B\in\Ccat^{\Osh}$). Para
$A\in\Ccat^{\Osh}$ fixo, definem-se funtores $h_A$ e $\hsh_A$ com valores, respectivamente,
na categoria $\Ccat$ dos grupos abelianos e na categoria
$\Ccat^{\Zsh}$ dos feixes abelianos sobre $X$, mediante as fórmulas:
\begin{equation}\tag{1.1}
  h_A(B)=\operatorname{Hom}_{\Osh}(X;A,B),
  \qquad
  \hsh_A(B)=\Homsh_{\Osh}(A,B).
\end{equation}

Os funtores $h_A$ e $\hsh_A$ são funtores covariantes exatos pela esquerda,
dos quais consideramos os funtores derivados à direita, denotados respectivamente por
$\operatorname{Ext}_{\Osh}^{p}(X;A,B)$ e $\Extsh_{\Osh}^{p}(A,B)$. Portanto,
por definição,
\begin{equation}\tag{1.2}
\left\{
\begin{aligned}
\operatorname{Ext}_{\Osh}^{p}(X;A,B)
  &= (R^{p}h_A)(B)
   = H^{p}\bigl(\operatorname{Hom}_{\Osh}(X;A,C(B))\bigr),\\
\Extsh_{\Osh}^{p}(A,B)
  &= (R^{p}\hsh_A)(B)
   = H^{p}\bigl(\Homsh_{\Osh}(A,C(B))\bigr),
\end{aligned}
\right.
\end{equation}

onde o símbolo $R^{p}$ designa a passagem aos funtores derivados à direita, e onde
$C(B)$ designa uma resolução injetiva arbitrária de $B$ em $\Ccat^{\Osh}$.
Denotemos por
\[
  \Gamma:\Ccat^{\Zsh}\longrightarrow\Ccat
\]
o funtor \fquote{seções}. Lembremos que seus funtores derivados à direita são
denotados por
\[
  B\longmapsto H^{p}(X,B).
\]

\src{170}
\begin{equation}\tag{1.3}
  H^{p}(X,B)=(R^{p}\Gamma)(B)=H^{p}\bigl(\Gamma(C(B))\bigr).
\end{equation}
Evidentemente, tem-se $h_A=\Gamma\hsh_A$; por outro lado, verifica-se que $\hsh_A$
transforma objetos injetivos em objetos acíclicos para $\Gamma$. Conclui-se do
modo habitual:

\result{PROPOSIÇÃO}{1} Para todo $\Osh$-Módulo $A$ existe um funtor
espectral cohomológico sobre $\Ccat^{\Osh}$, que converge ao funtor graduado
$\bigl(\operatorname{Ext}_{\Osh}^{*}(X;A,B)\bigr)$ e cujo termo inicial é
\begin{equation}\tag{1.4}
  E_{2}^{p,q}(A,B)=H^{p}\bigl(X,\Extsh_{\Osh}^{q}(A,B)\bigr).
\end{equation}
Dela se deduzem \fquote{homomorfismos de bordo} e uma sequência exata de cinco termos
que não escreveremos.

\result{COROLÁRIO}{1} Se $A$ é localmente isomorfo a $\Osh^{n}$, então têm-se
isomorfismos canônicos
\begin{equation}\tag{1.5}
  \operatorname{Ext}_{\Osh}^{p}(X;A,B)
  \overset{\sim}{\longleftarrow}
  H^{p}\bigl(X,\Homsh_{\Osh}(A,B)\bigr)
\end{equation}
(dados por homomorfismos de bordo da sequência espectral). Em particular, tem-se
um isomorfismo canônico
\begin{equation}\tag{1.6}
  \operatorname{Ext}_{\Osh}^{p}(X;\Osh,B)=H^{p}(X,B).
\end{equation}

Para utilizar estes resultados, é preciso determinar explicitamente os
$\Extsh_{\Osh}^{p}(A,B)$. São funtores que se calculam localmente,
isto é, se $U$ é um aberto de $X$, tem-se
\[
  \Extsh_{\Osh}^{p}(A,B)|U
  =\Extsh_{\Osh|U}^{p}(A|U,B|U),
\]
como resulta do fato de que a restrição a $U$ de um $\Osh$-Módulo injetivo
é um $(\Osh|U)$-Módulo injetivo. Além disso, para $x\in X$ fixo, há homomorfismos funtoriais
\begin{equation}\tag{1.7}
  \Homsh_{\Osh}(A,B)_x
  \longrightarrow
  \operatorname{Hom}_{\Osh_x}(A_x,B_x),
\end{equation}
que se prolongam de maneira única a um homomorfismo de funtores cohomológicos
em $B$:
\begin{equation}\tag{1.8}
  \Extsh_{\Osh}^{p}(A,B)_x
  \longrightarrow
  \operatorname{Ext}_{\Osh_x}^{p}(A_x,B_x).
\end{equation}

\result{PROPOSIÇÃO}{2} Se $A$ é, em uma vizinhança de $x$, isomorfo ao
conúcleo de um homomorfismo $\Osh^{m}\longrightarrow\Osh^{n}$, então (1.8) é um
isomorfismo para todo $p$.\footnote{O texto impresso diz « (1.7) »; a
presença de $p$ e a frase anterior exigem « (1.8) ».} Isto ocorre, em
particular, se $A$ é um $\Osh$-módulo coerente [3].

\src{171}
\result{PROPOSIÇÃO}{3} Seja $L_{*}=(L_i)$ uma resolução pela esquerda do
$\Osh$-Módulo $A$ por $\Osh$-Módulos cada um dos quais é localmente isomorfo a
um $\Osh^{n}$. Então $\Extsh_{\Osh}^{*}(A,B)$ identifica-se com
$H^{*}\bigl(\Homsh_{\Osh}(L_{*},B)\bigr)$, e
$\operatorname{Ext}_{\Osh}^{*}(X;A,B)$ identifica-se com a hipercohomologia de $X$
com respeito ao complexo $\Homsh_{\Osh}(L_{*},B)$.

A demonstração é padrão: considera-se o bicomplexo
$\Homsh_{\Osh}(L_{*},C(B))$, onde $C(B)$ é uma resolução injetiva de $B$, e os
homomorfismos naturais de $\Homsh_{\Osh}(L_{*},B)$ e de
$\Homsh_{\Osh}(A,C(B))$ neste último.

Para terminar, assinalemos que os dois funtores $\operatorname{Ext}$ introduzidos
em (1.2) não são apenas funtores cohomológicos em $B$, mas também
bifuntores cohomológicos, covariantes em $B$ e contravariantes em $A$.

\fsection{2}{Lei de composição nos Ext}

Os resultados deste número devem-se, independentemente, a CARTIER e a YONEDA; veja-se
uma exposição de CARTIER [1] para mais detalhes. Seja $\Ccat$ uma categoria abeliana.
Sejam $K$ e $L$ dois objetos graduados de $\Ccat$. Denotamos por
$\operatorname{Hom}(K,L)$ o grupo abeliano graduado cuja componente de grau
$n$ é formada pelos homomorfismos homogêneos de grau $n$ de $K$ em $L$ (isto é,
os sistemas $(u_i)$ de homomorfismos $K^{i}\longrightarrow L^{i+n}$).

Se $K$ e $L$ são complexos (com operadores diferenciais de grau $+1$, para
fixar as ideias), introduz-se em $\operatorname{Hom}(K,L)$ o operador
diferencial dado por
\begin{equation}\tag{2.1}
  \delta(u)=du+(-1)^{n+1}ud,\qquad\text{onde }n=\deg(u),
\end{equation}
que o converte em um complexo com operador diferencial de grau $+1$. Os ciclos de
grau $n$ são os operadores de grau $n$ que anticomutam com $d$ (como
operadores homogêneos).\footnote{O texto impresso diz « com $u$ »; a equação
(2.1) exige « com $d$ ».} Pode-se então considerar
$H^{*}\bigl(\operatorname{Hom}(K,L)\bigr)$; este é um invariante dos tipos
de homotopia de $K$ e de $L$, que poderemos denotar por $H^{*}(K,L)$.

Dado um terceiro complexo $M$, a composição dos homomorfismos
define um acoplamento
\[
  \operatorname{Hom}(K,L)\times\operatorname{Hom}(L,M)
  \longrightarrow\operatorname{Hom}(K,M)
\]
compatível com os operadores diferenciais, de onde, ao passar à cohomologia,
obtêm-se acoplamentos
\begin{equation}\tag{2.2}
  H^{*}(K,L)\times H^{*}(L,M)\longrightarrow H^{*}(K,M),
\end{equation}
denotados por $(u,v)\longmapsto vu$. Estes acoplamentos satisfazem a evidente propriedade
de associatividade. Em particular, $H^{*}(K,K)$ é um anel graduado associativo
com unidade, e $H^{*}(K,L)$ (resp. $H^{*}(L,K)$) é um módulo graduado pela direita
(resp. pela esquerda) sobre este anel, etc. Em grau $0$, (2.2) reduz-se à
composição dos morfismos de complexos. Finalmente, uma sequência exata de
complexos
\[
  0\longrightarrow K'\longrightarrow K\longrightarrow K''\longrightarrow 0,
\]
\src{172}
tal que $K'^{i}$ se identifica com um somando direto de $K^{i}$ para todo $i$,
dá lugar a uma sequência exata de complexos de grupos
$\operatorname{Hom}(K'',L)$, etc., e daí a um operador de conexão
\[
  H^{i}(K',L)\longrightarrow H^{i+1}(K'',L).
\]
Definem-se de maneira análoga os operadores de conexão relativos a uma sequência exata em $L$. Os
acoplamentos (2.2) são compatíveis, no sentido usual, com estes operadores de conexão.

Suponhamos agora que $\Ccat$ é uma categoria tal que todo objeto $A$
de $\Ccat$ admite uma resolução injetiva $C(A)$. Verifica-se então, utilizando
uma das numerosas variantes do teorema do bicomplexo, que
\[
  H^{*}\bigl(C(A),C(B)\bigr)
  =H^{*}\bigl(\operatorname{Hom}(C(A),C(B))\bigr)
\]
é canonicamente isomorfo a
\[
  H^{*}\bigl(\operatorname{Hom}(A,C(B))\bigr)
  =\operatorname{Ext}^{*}(A,B).
\]
Os operadores de conexão considerados anteriormente dão os homomorfismos de conexão dos
$\operatorname{Ext}$. Além disso, os acoplamentos (2.2) dão aqui
acoplamentos associativos, compatíveis com os homomorfismos de conexão:
\begin{equation}\tag{2.3}
  \operatorname{Ext}^{*}(A,B)\times\operatorname{Ext}^{*}(B,C)
  \longrightarrow\operatorname{Ext}^{*}(A,C).
\end{equation}
Em particular, $\operatorname{Ext}^{*}(A,A)$ é um anel graduado associativo
com unidade, etc. (Demonstra-se de maneira análoga que os funtores
$\operatorname{Ext}$ operam nos funtores derivados de um funtor
qualquer; não utilizaremos aqui este fato.)

No caso em que a categoria considerada seja a categoria $\Ccat^{\Osh}$ dos
$\Osh$-Módulos sobre $X$, obtêm-se, portanto, acoplamentos
\begin{equation}\tag{2.4}
  \operatorname{Ext}_{\Osh}^{p}(X;A,B)
  \times\operatorname{Ext}_{\Osh}^{q}(X;B,C)
  \longrightarrow\operatorname{Ext}_{\Osh}^{p+q}(X;A,C),
\end{equation}
que podem ser explicitados como foi indicado. De resto, os mesmos
desenvolvimentos, substituindo a categoria dos grupos abelianos pela
categoria dos feixes abelianos sobre $X$, e os funtores $\operatorname{Hom}$
pelos funtores $\Homsh$, definem também acoplamentos com as mesmas
propriedades formais e, desta vez, de \fquote{natureza local}:
\begin{equation}\tag{2.5}
  \Extsh_{\Osh}^{p}(A,B)\times\Extsh_{\Osh}^{q}(B,C)
  \longrightarrow\Extsh_{\Osh}^{p+q}(A,C).
\end{equation}
Estes últimos são determinados observando-se que os homomorfismos (1.8) são
compatíveis com os acoplamentos entre os $\operatorname{Ext}$.

\src{173}
Lembremos, por fim, que existe também uma estrutura multiplicativa entre os funtores
$H^{p}(X,A)$ (produto cup). Verifica-se então que as sequências espectrais da
proposição 1 são compatíveis com as estruturas multiplicativas; mais
precisamente, tem-se um acoplamento da sequência espectral $E(A,B)$ e da sequência
espectral $E(B,C)$ para a sequência espectral $E(A,C)$ que, no limite,
se reduz ao acoplamento entre os $\operatorname{Ext}$ globais e, no
termo inicial, ao acoplamento deduzido, no segundo membro de (1.4), do
produto cup e dos acoplamentos dos $\operatorname{Ext}$ locais. Disso resulta,
em particular, que os \fquote{homomorfismos de bordo}
\begin{equation}\tag{2.6}
  \operatorname{Ext}_{\Osh}^{n}(X;A,B)
  \longrightarrow H^{0}\bigl(X,\Extsh_{\Osh}^{n}(A,B)\bigr)
\end{equation}
e
\begin{equation}\tag{2.7}
  H^{n}\bigl(X,\Homsh_{\Osh}(A,B)\bigr)
  \longrightarrow\operatorname{Ext}_{\Osh}^{n}(X;A,B)
\end{equation}
são compatíveis com as estruturas multiplicativas. Portanto, se nos restringirmos aos
feixes localmente isomorfos a um $\Osh^{m}$, isto explicita por completo a
composição entre os $\operatorname{Ext}$ globais mediante o produto cup, levando
em conta o isomorfismo (1.5).

\fsection{3}{Resultados de cohomologia local}

Seja $A$ um anel comutativo com unidade provido de um ideal $J$. Vamos
definir, para um $A$-módulo variável $M$, homomorfismos funtoriais
\begin{equation}\tag{3.1}
\left\{
\begin{aligned}
\operatorname{Ext}_{A}^{p}(A/J,M)
  &\longrightarrow
  \operatorname{Hom}_{A}\!\left(\bigwedge^{p}J/J^{2},M\otimes A/J\right),\\
\operatorname{Tor}_{p}^{A}(A/J,M)
  &\longleftarrow
  \left(\bigwedge^{p}J/J^{2}\right)
  \otimes\operatorname{Hom}_{A}(A/J,M),
\end{aligned}
\right.
\end{equation}
onde os produtos tensoriais e exteriores são tomados sobre o anel $A$; note-se,
por outro lado, que $J/J^{2}$ é, na realidade, um $A/J$-módulo, e que suas potências
exteriores como $A$-módulo ou como $A/J$-módulo são as mesmas.
Definir os homomorfismos (3.1) equivale a definir, para todo
sistema $\underline{x}=(x_1,\ldots,x_p)$ de elementos de $J$, homomorfismos
$\varphi_{\underline{x}}$:
\begin{equation}\tag{3.2}
\left\{
\begin{aligned}
\varphi_{\underline{x}}:
  \operatorname{Ext}_{A}^{p}(A/J,M)&\longrightarrow M\otimes A/J,\\
\varphi_{\underline{x}}:
  \operatorname{Hom}_{A}(A/J,M)&\longrightarrow
  \operatorname{Tor}_{p}^{A}(A/J,M),
\end{aligned}
\right.
\end{equation}
de tal maneira que sejam satisfeitas as condições seguintes:

\src{174}
\begin{enumerate}
\item[\textup{i.}] $\varphi_{x_1,\ldots,x_p}$ depende de maneira
$A$-multilinear alternada do sistema dos $x_i\in J$.
\item[\textup{ii.}] $\varphi_{x_1,\ldots,x_p}$ é nulo quando algum dos $x_i$ está
em $J^{2}$.
\end{enumerate}
Na realidade, ii. resulta de i., pois $a\varphi_{\underline{x}}=0$ para $a\in J$,
como se vê observando que todos os módulos de (3.2) são anulados por
$J$.

Para definir $\varphi_{\underline{x}}$, consideremos o complexo
$K_{\underline{x}}$ cujo $A$-módulo subjacente é $\bigwedge A^{p}$, e
cujo operador diferencial é o produto interior $i_{\underline{x}}$ por
$\underline{x}$ considerado como forma linear sobre $A^{p}$ de componentes
$x_1,\ldots,x_p$. O operador diferencial é de grau $-1$, os graus são
positivos, e a cohomologia deste complexo em grau $0$ é
$A/(x_1A+\cdots+x_pA)$. Como os $x_i$ estão em $J$, deduz-se uma
aumentação $K_{\underline{x},0}\longrightarrow A/J$. Assim,
$K_{\underline{x}}$ aparece como um complexo livre aumentado, com módulo
de aumentação $A/J$. Disso se deduzem, do modo habitual, homomorfismos
\[
\operatorname{Ext}_{A}^{*}\bigl(H_0(K_{\underline{x}}),M\bigr)
\longrightarrow
H^{*}\bigl(\operatorname{Hom}_{A}(K_{\underline{x}},M)\bigr)
\]
e
\[
\operatorname{Tor}_{*}^{A}\bigl(H_0(K_{\underline{x}}),M\bigr)
\longleftarrow
H_{*}\bigl(K_{\underline{x}}\otimes_A M\bigr),
\]
de onde, compondo com os homomorfismos sobre os $\operatorname{Ext}$ e os
$\operatorname{Tor}$ deduzidos do homomorfismo de aumentação
$H_0(K_{\underline{x}})\longrightarrow A/J$, obtêm-se homomorfismos
\begin{equation}\tag{3.3}
\left\{
\begin{aligned}
\psi_{\underline{x}}:
\operatorname{Ext}_{A}^{*}(A/J,M)
&\longrightarrow H^{*}\bigl(\operatorname{Hom}_{A}(K_{\underline{x}},M)\bigr),\\
\psi_{\underline{x}}:
\operatorname{Tor}_{*}^{A}(A/J,M)
&\longleftarrow H_{*}\bigl(K_{\underline{x}}\otimes_A M\bigr).
\end{aligned}
\right.
\end{equation}
Ora, verifica-se imediatamente que, em grau máximo $p$, a cohomologia dos segundos
membros é $M/(x_1M+\cdots+x_pM)$ (resp. é o conjunto dos elementos de $M$
anulados por cada um dos $x_i$). Como os $x_i$ estão em $J$, deduzem-se
homomorfismos
\begin{equation}\tag{3.4}
\left\{
\begin{aligned}
H^{p}\bigl(\operatorname{Hom}_{A}(K_{\underline{x}},M)\bigr)
&\longrightarrow M\otimes A/J,\\
H_{p}\bigl(K_{\underline{x}}\otimes_A M\bigr)
&\longleftarrow\operatorname{Hom}_{A}(A/J,M).
\end{aligned}
\right.
\end{equation}
Compondo os homomorfismos (3.3) e (3.4), obtêm-se os homomorfismos (3.2)
que queríamos definir. A verificação de i. é laboriosa, mas não
apresenta dificuldades.

\src{175}
\result{PROPOSIÇÃO}{4} Seja $A$ um anel comutativo com unidade, e seja
$(x_1,\ldots,x_p)$ uma sequência de elementos de $A$ tal que, para $1\leq i\leq p$,
a imagem de $x_i$ no quociente de $A$ pelo ideal gerado por
$(x_1,\ldots,x_{i-1})$ não seja divisor de $0$. Seja $J$ o ideal gerado
pelos $x_i$. Então $J/J^{2}$ é um $(A/J)$-módulo livre que tem por base as
imagens canônicas dos $x_i$, o complexo $K_{\underline{x}}$ é uma resolução
livre de $A/J$, e para todo $A$-módulo $M$, os homomorfismos (3.1) em
grau $p$ são bijetivos. O mesmo ocorre com os homomorfismos análogos
definidos para graus $i$ quaisquer, sempre que $J.M=0$.

(O ponto essencial, do qual resultam todos os demais, é a aciclicidade de
$K_{\underline{x}}$, que é um fato bem conhecido sob as condições indicadas.)

\result{COROLÁRIO}{1} Sejam $A$ e $J$ como acima, e suponhamos além disso que
$A$ é uma álgebra afim regular de $\dim n$ sobre um corpo perfeito $k$, e
que $A/J$ é uma álgebra afim regular. Denotemos por $\Omega^{1}(A)$,
$\Omega^{1}(A/J)$ os módulos de diferenciais de Kähler. Então tem-se um
isomorfismo canônico compatível com a localização,
\begin{equation}\tag{3.5}
  \operatorname{Ext}_{A}^{p}\bigl(\Omega^{n-p}(A/J),\Omega^{n}(A)\bigr)=A/J.
\end{equation}
Com efeito, $\Omega^{n-p}(A/J)$ é um $(A/J)$-módulo livre de posto $1$, e igualmente
$\Omega^{n}(A)$ é um $A$-módulo livre de posto $1$,\footnote{O texto impresso
diz « posto $n$ »; os dois módulos de formas de grau máximo têm posto
$1$, como exige também a identificação com $\bigwedge^{p}(J/J^{2})$ mais abaixo.}
pelo que o primeiro membro é igual a
\[
  \operatorname{Ext}_{A}^{p}(A/J,A)
  \otimes\Omega^{n-p}(A/J)'\otimes\Omega^{n}(A)
\]
(onde o símbolo « $'$ » designa o $(A/J)$-módulo dual). O produto tensorial dos
dois últimos fatores identifica-se com $\bigwedge^{p}(J/J^{2})$, pelo que o conjunto
identifica-se com
\[
  \operatorname{Ext}_{A}^{p}\!\left(A/J,\bigwedge^{p}(J/J^{2})\right),
\]
e, pela proposição, com
\[
  \operatorname{Hom}_{A}\!\left(\bigwedge^{p}J/J^{2},
  \bigwedge^{p}J/J^{2}\right),
\]
isto é, com $A/J$. Em particular, em
$\operatorname{Ext}_{A}^{p}\bigl(\Omega^{n-p}(A/J),\Omega^{n}(A)\bigr)$ há um
elemento distinguido, correspondente ao elemento unidade de $A/J$, chamado
\emph{classe fundamental do ideal $J$ em $A$}. (Na realidade, pode
ser definida sob condições sensivelmente mais gerais.) O
corolário 1 pode ser apresentado em uma forma mais geométrica e mais global:

\result{COROLÁRIO}{2} Sejam $X$ uma variedade não singular sobre um corpo
algebricamente fechado $k$, $Y$ uma subvariedade fechada não singular de $X$,
$\Osh_X$ o feixe de anéis locais de $X$, e $\Osh_Y$ o feixe de anéis
locais de $Y$ considerado como feixe quociente de $\Osh_X$. Sejam $n$ a
dimensão de $X$ e $n-p$ a de $Y$.

\src{176}
Seja $\Omega_X$ (resp. $\Omega_Y$) o feixe de germes de formas
diferenciais regulares sobre $X$ (resp. $Y$). Então têm-se isomorfismos
canônicos:
\begin{equation}\tag{3.6}
  \Extsh_{\Osh_X}^{p}\bigl(\Omega_Y^{n-p},\Omega_X^{n}\bigr)=\Osh_Y,
\end{equation}
ou também
\begin{equation}\tag{3.6 bis}
  \Extsh_{\Osh_X}^{p}\bigl(\Osh_Y,\Omega_X^{n}\bigr)=\Omega_Y^{n-p}.
\end{equation}
A fórmula (3.6 bis) pode servir como definição de $\Omega_Y^{n-p}$ quando $Y$
é uma variedade singular. Mais precisamente:

\result{PROPOSIÇÃO}{5} Seja $Y$ um subconjunto algébrico de dimensão
$q=n-p$ em uma variedade algébrica não singular $X$ de dimensão $n$. Sejam
$F$ um feixe algébrico coerente sobre $X$ cujo suporte está contido em $Y$, e $L$
um feixe algébrico localmente livre sobre $X$. Então os feixes
$\Extsh_{\Osh_X}^{i}(F,L)$ são nulos para $i<p$, enquanto que para $i=p$ tem-se um
isomorfismo canônico
\begin{equation}\tag{3.7}
  \Extsh_{\Osh_X}^{p}(F,L)
  =\Homsh_{\Osh_X}\!\left(F,
  \Extsh_{\Osh_X}^{p}(\Osh_X/J,L)\right),
\end{equation}
onde $J$ designa um feixe de ideais arbitrário sobre $X$ que anula $F$ e cujo
conjunto de zeros é $Y$. Em particular, se $F$ é um feixe algébrico
coerente sobre $Y$, tem-se
\begin{equation}\tag{3.7 bis}
  \Extsh_{\Osh_X}^{p}(F,L)
  =\Homsh_{\Osh_Y}\!\left(F,\Extsh_{\Osh_X}^{p}(\Osh_Y,L)\right).
\end{equation}
Por fim, sendo ainda $F$ um feixe algébrico coerente sobre $Y$, os feixes
\[
  E^{i}(F)=\Extsh_{\Osh_X}^{p+i}(F,\Omega_X^{n})
\]
não dependem da maneira de mergulhar o espaço algébrico $Y$ em uma variedade
não singular $X$.

Como a questão é local, pode-se supor que $X$ é afim e $L=\Osh_X$. Isto reduz
o problema a uma questão de álgebra comutativa, e até de álgebra local: se $A$ é um
anel local regular e $M$ um $A$-módulo cujo suporte tem dimensão
$\leq q=n-p$, é preciso demonstrar que $\operatorname{Ext}_{A}^{i}(M,A)=0$ para $i<p$, e
\[
  \operatorname{Ext}_{A}^{p}(M,A)
  =\operatorname{Hom}_{A}\!\left(M,\operatorname{Ext}_{A}^{p}(A/J,A)\right),
\]
onde $J$ é um ideal qualquer de « dimensão » $\leq q$ que anula $M$. Para o
primeiro ponto, procede-se por indução sobre $q$: um dévissage imediato reduz
ao caso em que $M$ é da forma $A/J$ e, em seguida, substituindo $J$ por um ideal menor
e utilizando a hipótese de indução e a sequência exata dos
$\operatorname{Ext}$, ao caso em que $J$ é gerado por um « sistema de parâmetros »
como na proposição 4, onde o resultado é imediato. O resultado
anterior implica que, se $J$ é um ideal

\src{177}
fixo de « dimensão » $\leq q$, então o funtor contravariante
$E(M)=\operatorname{Ext}_{A}^{p}(M,A)$ sobre a categoria dos $A/J$-módulos é
exato pela esquerda; além disso, transforma somas diretas em produtos diretos,
de onde resulta facilmente que
\[
  E(M)=\operatorname{Hom}_{A}\bigl(M,E(A/J)\bigr).
\]
\footnote{O texto impresso diz $E(A)$; na categoria considerada, o objeto
representante é $A/J$, e a igualdade requerida para (3.7) é $E(A/J)$.} Por fim, a
última afirmação da proposição 5 é mais sutil e resulta de uma
caracterização intrínseca dos $E^{i}(F)$ mediante um teorema de dualidade
local que não pode ser exposto aqui.

\result{COROLÁRIO}{} Denotemos por $\omega_Y^{q}$ o feixe
$\Extsh_{\Osh_X}^{p}(\Osh_Y,\Omega_X^{n})$. Então tem-se um isomorfismo
funtorial para os feixes algébricos coerentes $F$ sobre $Y$:
\begin{equation}\tag{3.8}
  \Extsh_{\Osh_X}^{p}(F,\Omega_X^{n})
  =\Homsh_{\Osh_X}(F,\omega_Y^{q}).
\end{equation}

\fsection{4}{Classe de cohomologia associada a uma subvariedade}

Em tudo o que segue, $X$ designa um conjunto algébrico de dimensão $n$, definido
sobre um corpo $k$, que suporemos algebricamente fechado para simplificar. Salvo
no n\textsuperscript{o} 6, supõe-se que $X$ é não singular. Denotamos por
$\Osh_X$ o feixe de anéis locais sobre $X$, e por
\[
  \Omega_X^{*}=\bigoplus_p\Omega_X^{p}
\]
o feixe de germes de formas diferenciais sobre $X$.\footnote{O texto
impresso emprega um sinal de união; a estrutura graduada é a soma direta
dos feixes $\Omega_X^p$.} Se $Y$ é uma parte fechada de $X$, identificamos
os feixes algébricos coerentes sobre $Y$ com feixes coerentes sobre $X$
nulos fora de $Y$; em particular, assim ocorre com $\Osh_Y$ e $\Omega_Y$.

\result{LEMA}{1} Seja $F$ um feixe algébrico coerente sobre $X$ cujo
suporte tem dimensão $\leq n-p$, e $L$ um feixe algébrico coerente
localmente livre sobre $X$. Então $\operatorname{Ext}_{\Osh_X}^{i}(X;F,L)$ é nulo
para $i<p$, e tem-se um isomorfismo canônico
\begin{equation}\tag{4.1}
  \operatorname{Ext}_{\Osh_X}^{p}(X;F,L)
  =H^{0}\bigl(X,\Extsh_{\Osh_X}^{p}(F,L)\bigr).
\end{equation}
Se $F$ é um feixe algébrico coerente sobre uma parte fechada $Y$ de $X$ de
dimensão $\leq n-p$, tem-se um isomorfismo canônico
\begin{equation}\tag{4.1 bis}
  \Extsh_{\Osh_X}^{p}(F,L)
  =\Homsh_{\Osh_X}\!\left(
    F\otimes L'\otimes\Omega_X^{n},\omega_Y^{n-p}
  \right),
\end{equation}
onde $\omega_Y^{n-p}$ é o feixe sobre $Y$ definido no corolário da
proposição 5 (que se identifica com $\Omega_Y^{n-p}$ se $Y$ é não singular).

A fórmula (4.1) é uma consequência imediata da sequência espectral da
proposição 1 e da proposição 5; em virtude da fórmula (3.8), pode-se
escrever

\src{178}
\begin{align*}
\Extsh_{\Osh_X}^{p}(F,L)
&=L\otimes(\Omega_X^{n})'\otimes
  \Extsh_{\Osh_X}^{p}(F,\Omega_X^{n})\\
&=L\otimes(\Omega_X^{n})'\otimes
  \Homsh_{\Osh_X}(F,\omega_Y^{q})\\
&=\Homsh_{\Osh_X}\!\left(
  F\otimes L'\otimes\Omega_X^{n},\omega_Y^{q}\right),
\end{align*}
onde $q=n-p$, de onde resulta a fórmula (4.1 bis).

Tomando, em particular, $F=\Osh_Y$, $L=\Omega_X^{p}$, obtém-se, tendo em conta
que $\Omega_X^{n}\otimes(\Omega_X^{p})'=\Omega_X^{n-p}$, um isomorfismo
canônico
\begin{equation}\tag{4.2}
  \operatorname{Ext}_{\Osh_X}^{p}(X;\Osh_Y,\Omega_X^{p})
  =\operatorname{Hom}_{\Osh_X}
  (X;\Omega_X^{n-p},\omega_Y^{n-p}).
\end{equation}
Suponhamos agora que $Y$ é não singular, para simplificar, de modo que
$\omega_Y^{n-p}=\Omega_Y^{n-p}$. Tem-se um homomorfismo natural de
$\Omega_X^{n-p}$ em $\Omega_Y^{n-p}$, de onde resulta uma seção canônica $s_Y$ do
feixe $\Extsh_{\Osh_X}^{p}(\Osh_Y,\Omega_X^{p})$, que chamaremos, se
todas as componentes de $Y$ têm dimensão $n-p$, \emph{seção
fundamental} do feixe $\Extsh_{\Osh_X}^{p}(\Osh_Y,\Omega_X^{p})$. Esta
seção define, em virtude de (4.1), um elemento de
$\operatorname{Ext}_{\Osh_X}^{p}(X;\Osh_Y,\Omega_X^{p})$. Ora, o homomorfismo
natural $\Osh_X\longrightarrow\Osh_Y$ define um homomorfismo
\[
  \operatorname{Ext}_{\Osh_X}^{p}(X;\Osh_Y,\Omega_X^{p})
  \longrightarrow
  \operatorname{Ext}_{\Osh_X}^{p}(X;\Osh_X,\Omega_X^{p})
  =H^{p}(X,\Omega_X^{p}).
\]
Obtém-se assim um elemento de $H^{p}(X,\Omega_X^{p})$, denotado por $P_X(Y)$ e chamado
\emph{classe de cohomologia de $Y$ em $X$}. Deduz-se, portanto, da seção
$s_Y$ de $\Extsh_{\Osh_X}^{p}(\Osh_Y,\Omega_X^{p})$ mediante o diagrama
de homomorfismos seguinte:
\begin{equation}\tag{4.3}
\begin{aligned}
H^{p}(X,\Omega_X^{p})
&=\operatorname{Ext}_{\Osh_X}^{p}(X;\Osh_X,\Omega_X^{p})\\
&\longleftarrow\operatorname{Ext}_{\Osh_X}^{p}(X;\Osh_Y,\Omega_X^{p})
\xrightarrow{\ \sim\ }
H^{0}\!\left(X,\Extsh_{\Osh_X}^{p}(\Osh_Y,\Omega_X^{p})\right).
\end{aligned}
\end{equation}

Chamemos \emph{ciclo não singular de dimensão $n-p$} a um elemento do grupo
abeliano livre gerado pelas subvariedades irredutíveis não singulares de
dimensão $n-p$ em $X$. Então a função $Y\longmapsto P(Y)$ se prolonga a
um homomorfismo do grupo de ciclos não singulares de dimensão $n-p$ sobre $X$
no grupo $H^{p}(X,\Omega_X^{p})$.

Sejam $Z^{n-p}$ e $Z'^{n-p'}$ ciclos não singulares de dimensões $n-p$ e
$n-p'$; diz-se que se \emph{intersectam transversalmente} se toda componente
de $Z$ interseca transversalmente toda componente de $Z'$. Então o ciclo $Z.Z'$ está
definido e é um ciclo não singular de dimensão $n-p-p'$. Dito isto, tem-se o

\src{179}
\result{TEOREMA}{1} Se $Z^{n-p}$ e $Z'^{n-p'}$ são ciclos não singulares
que se intersectam transversalmente, então tem-se
\begin{equation}\tag{4.4}
  P_X(Z.Z')=P_X(Z).P_X(Z'),
\end{equation}
onde o produto do segundo membro é o produto cup
\[
  H^{p}(X,\Omega_X^{p})\times H^{p'}(X,\Omega_X^{p'})
  \longrightarrow H^{p+p'}(X,\Omega_X^{p+p'}).
\]
(Supõe-se que $X$ é isomorfo a uma parte localmente fechada de um espaço
projetivo.)

Esta última hipótese serve apenas para poder concluir que todo
feixe algébrico coerente sobre $X$ é quociente de um feixe algébrico
coerente localmente livre (SERRE) e, portanto, admite uma resolução pela esquerda mediante
feixes localmente livres. Para demonstrar o teorema 1, pode-se supor que
$Z$ e $Z'$ são subvariedades irredutíveis não singulares $Y$ e $Y'$ que se
intersectam transversalmente. Seja $L_{*}$ uma resolução pela esquerda de $\Osh_Y$ mediante
feixes localmente livres; então, em virtude da proposição 3, o diagrama
de homomorfismos (4.3) identifica-se com o diagrama
\[
R^{p}\Gamma(\Omega_X^{p})
\xleftarrow{\alpha}
\mathbb R^{p}\Gamma\!\left(\Homsh_{\Osh_X}(L_{*},\Omega_X^{p})\right)
\xrightarrow[\sim]{\beta}
\Gamma\!\left(H^{p}\bigl(\Homsh_{\Osh_X}(L_{*},\Omega_X^{p})\bigr)\right),
\]
onde $\Gamma$ é o funtor « grupo de seções » sobre a categoria dos
feixes abelianos sobre $X$, e $\mathbb R^{p}\Gamma$ designa
a hipercohomologia de grau $p$ desse funtor, e $R^{p}\Gamma$ seu
$p$-ésimo funtor derivado. Suporemos, para simplificar, que $L_0=\Osh_X$ e que
a aumentação $L_0\longrightarrow\Osh_Y$ é o homomorfismo natural (o que
é lícito); então $\alpha$ se deduz do homomorfismo de complexos
$\Osh_X\longrightarrow L_{*}$ (considerando $\Osh_X$ como um complexo concentrado
em grau $0$), tendo em conta que $\mathbb R^{p}\Gamma(K)=R^{p}\Gamma(K_0)$ se $K$
é um complexo de feixes concentrado em grau $0$. O homomorfismo $\beta$ é um
« homomorfismo de bordo » bem conhecido. Consideremos um diagrama análogo, relativo a
uma resolução localmente livre $L'_{*}$ de $\Osh_{Y'}$, e consideremos o
diagrama comutativo de acoplamentos:
\begin{equation}\tag{4.5}
\resizebox{\textwidth}{!}{$
\begin{array}{ccc}
R^{p}\Gamma(\Omega_X^{p})
&\longleftarrow
\mathbb R^{p}\Gamma\!\left(\Homsh_{\Osh_X}(L_{*},\Omega_X^{p})\right)
&\xrightarrow{\sim}
\Gamma H^{p}\!\left(\Homsh_{\Osh_X}(L_{*},\Omega_X^{p})\right)
\\[3pt]
\times&\times&\times\\[3pt]
R^{p'}\Gamma(\Omega_X^{p'})
&\longleftarrow
\mathbb R^{p'}\Gamma\!\left(\Homsh_{\Osh_X}(L'_{*},\Omega_X^{p'})\right)
&\xrightarrow{\sim}
\Gamma H^{p'}\!\left(\Homsh_{\Osh_X}(L'_{*},\Omega_X^{p'})\right)
\\[3pt]
\downarrow&\downarrow&\downarrow\\[3pt]
R^{p+p'}\Gamma(\Omega_X^{p+p'})
&\longleftarrow
\mathbb R^{p+p'}\Gamma\!\left(
\Homsh_{\Osh_X}(L_{*}\otimes L'_{*},\Omega_X^{p+p'})\right)
&\xrightarrow{\sim}
\Gamma H^{p+p'}\!\left(
\Homsh_{\Osh_X}(L_{*}\otimes L'_{*},\Omega_X^{p+p'})\right).
\end{array}
$}
\end{equation}

\src{180}
Os acoplamentos das duas colunas da direita se deduzem do acoplamento de
complexos de feixes
\[
\Homsh_{\Osh_X}(L_{*},\Omega_X^{p})
\times\Homsh_{\Osh_X}(L'_{*},\Omega_X^{p'})
\longrightarrow
\Homsh_{\Osh_X}(L_{*}\otimes L'_{*},\Omega_X^{p+p'}),
\]
que se define mediante o produto exterior
$\Omega_X^{p}\times\Omega_X^{p'}\longrightarrow\Omega_X^{p+p'}$.
O acoplamento da primeira coluna é o produto cup (relativo ao produto
exterior). Afirmo que a última linha de (4.5) identifica-se com o diagrama
de isomorfismos análogo a (4.3), onde $Y$ é substituído por $Y\cap Y'$ e $p$ por
$p+p'$. Para isso basta demonstrar que $L\otimes L'$ é uma resolução
(evidentemente localmente livre) de $\Osh_{Y\cap Y'}$. Com efeito, tem-se
\[
  H_0(L\otimes L')=\Osh_Y\otimes\Osh_{Y'}=\Osh_{Y\cap Y'}
\]
e
\[
  H_i(L\otimes L')=\operatorname{Tor}_{i}^{\Osh_X}(\Osh_Y,\Osh_{Y'})=0
  \qquad(i>0),
\]
pelo fato de que $Y$ e $Y'$ se intersectam transversalmente. O teorema 1 resulta
agora da fórmula:
\begin{equation}\tag{4.6}
  s_Y.s_{Y'}=s_{Y.Y'}
\end{equation}
(onde o produto do primeiro membro é o da última coluna de (4.5)).
Esta fórmula (4.6), de natureza puramente local, verifica-se sem dificuldade
tomando para $L_{*}$ e $L'_{*}$ as resoluções consideradas na proposição
4. Demonstra-se igualmente (mediante uma demonstração mais simples) que
$Z\longmapsto P_X(Z)$ é compatível com o produto cartesiano:
\begin{equation}\tag{4.7}
  P_{X\times X'}(Z\times Z')=P_X(Z)\otimes P_{X'}(Z')
\end{equation}
(fórmula válida se $Z$, resp. $Z'$, é um ciclo não singular sobre a variedade
não singular $X$, resp. $X'$, considerando $Z\times Z'$ como um ciclo não
singular sobre $X\times X'$). De (4.4) e (4.7) deduz-se que $P_X(Z)$ é também
compatível com a operação « imagem inversa » por um morfismo de variedades não
singulares $f:X\longrightarrow X'$:
\begin{equation}\tag{4.8}
  P_X\bigl(f^{-1}(Z')\bigr)=f^{*}\bigl(P_{X'}(Z')\bigr),
\end{equation}
fórmula válida se $Z'$ é um ciclo não singular sobre $X'$ tal que $f$ seja
« transversal » a $Z'$, isto é, tal que a gráfica de $f$ seja transversal ao ciclo
$X\times Z'$ em $X\times X'$.

\src{181}
\result{COROLÁRIO}{1} Sejam $X$, $X'$ duas variedades não singulares
localmente fechadas em um espaço projetivo, e suponhamos que $X'$ é completa. Sejam $U$
um ciclo não singular sobre $X\times X'$, $a$ e $b$ dois pontos de $X'$ tais que
$U$ interseca transversalmente os ciclos $X\times(a)$ e $X\times(b)$, e sejam $Z$
e $Z'$ os ciclos não singulares sobre $X$ tais que
\[
  Z\times(a)=(X\times(a)).U,
  \qquad
  Z'\times(b)=(X\times(b)).U.
\]
Então tem-se
\[
  P_X(Z)=P_X(Z').
\]

Com efeito, seja $f_a:X\longrightarrow X\times X'$ definido por
$f_a(x)=(x,a)$; então, em virtude de (4.8), tem-se a fórmula
$P(Z)=f_a^{*}(P(U))$, e analogamente $P(Z')=f_b^{*}(P(U))$. Ora, utilizando a fórmula de
Künneth
\[
  H^{*}(X\times X',\Omega_{X\times X'}^{*})
  =H^{*}(X,\Omega_X^{*})\otimes H^{*}(X',\Omega_{X'}^{*})
\]
e o fato de que $H^{0}(X',\Osh_{X'})$ se reduz aos escalares, obtém-se
facilmente $f_a^{*}=f_b^{*}$, de onde se segue a conclusão.

Para todo $x\in X$, $(x)$ é uma subvariedade não singular de $X$ de
codimensão $n$, e define portanto um elemento $\varepsilon_x$ de
$H^{n}(X,\Omega_X^{n})$. Se $X$ é uma variedade projetiva não singular, do
corolário 1 resulta que $\varepsilon_x$ não depende do ponto $x$ escolhido; denota-se
por $\varepsilon_X$ e chama-se \emph{classe fundamental} de
$H^{n}(X,\Omega_X^{n})$.

\result{OBSERVAÇÃO}{} Para dispor de uma teoria satisfatória, seria preciso definir
$P_X(Z)$ para ciclos $Z$ quaisquer e demonstrar o teorema 1 para uma
interseção própria de ciclos. (No momento de escrever esta exposição, isso ainda não foi
feito com toda generalidade.) Admitindo que tenha sido feito, o corolário 1
passa a ser: se $Z$ e $Z'$ são dois ciclos algebricamente equivalentes, então
$P_X(Z)=P_X(Z')$ (afirmação que não parece resultar do que precede, mesmo que
$Z$ e $Z'$ sejam não singulares).

\fsection{5}{O teorema de dualidade}

Neste número, $X$ designa uma variedade projetiva não singular de dimensão
$n$.

\result{TEOREMA}{2} A classe fundamental $\varepsilon_X$ de
$H^{n}(X,\Omega_X^{n})$ é uma base deste espaço vetorial.

(A demonstração será dada mais adiante.) Em virtude do teorema anterior, pode-se
identificar $H^{n}(X,\Omega_X^{n})$ com o corpo $k$. Consideremos agora os
acoplamentos considerados no n\textsuperscript{o} 2, que dão, em particular, um
acoplamento

\src{182}
\[
\operatorname{Ext}_{\Osh_X}^{p}(X;\Osh_X,F)
\times\operatorname{Ext}_{\Osh_X}^{n-p}(X;F,\Omega_X^{n})
\longrightarrow
\operatorname{Ext}_{\Osh_X}^{n}(X;\Osh_X,\Omega_X^{n}),
\]
isto é,
\begin{equation}\tag{5.1}
  H^{p}(X,F)
  \times\operatorname{Ext}_{\Osh_X}^{n-p}(X;F,\Omega_X^{n})
  \longrightarrow H^{n}(X,\Omega_X^{n}).
\end{equation}
Tendo em conta o teorema 2, este acoplamento define um homomorfismo
\begin{equation}\tag{5.2}
  \operatorname{Ext}_{\Osh_X}^{n-p}(X;F,\Omega_X^{n})
  \longrightarrow\bigl(H^{p}(X,F)\bigr)'.
\end{equation}
Este homomorfismo é funtorial em $F$ e comuta com os homomorfismos de conexão
relativos às sequências exatas $0\longrightarrow F'\longrightarrow F
\longrightarrow F''\longrightarrow0$.

\result{TEOREMA}{3} O homomorfismo (5.2) é um isomorfismo.

Recupera-se, em particular, o resultado de Serre:

\result{COROLÁRIO}{} Seja $E$ um fibrado vetorial algébrico sobre $X$ e
$\Osh_X(E)$ o feixe de germes de seções regulares de $E$; então têm-se
isomorfismos canônicos:
\begin{equation}\tag{5.3}
  \bigl(H^{p}(X,\Osh_X(E))\bigr)'
  =H^{n-p}\bigl(X,\Omega_X^{n}\otimes\Osh_X(E')\bigr).
\end{equation}
Basta aplicar o teorema 3 e o corolário 1 à proposição 1.

Os teoremas 2 e 3 resultarão do enunciado seguinte:

\noindent\textbf{(D) O homomorfismo seguinte}
\begin{equation}\tag{5.2 bis}
  \operatorname{Ext}_{\Osh_X}^{n-p}(X;F,\Omega_X^{n})
  \longrightarrow\bigl(H^{p}(X,F)\bigr)'\otimes L,
  \qquad L=H^{n}(X,\Omega_X^{n}),
\end{equation}
deduzido do acoplamento (5.1) é um isomorfismo.

Demonstremos, com efeito, que (D) implica o teorema 2. Seja
$k_x=\Osh_{(x)}$ o feixe de anéis locais da variedade reduzida ao ponto
$x\in X$; consideremos o homomorfismo canônico
$\Osh_X\longrightarrow k_x$ e o homomorfismo associado
\begin{equation}\tag{5.3 bis}
  H^{0}(X,\Osh_X)\longrightarrow H^{0}(X,k_x),
\end{equation}
seu transposto identifica-se com o homomorfismo
\begin{equation}\tag{5.4}
  \operatorname{Ext}_{\Osh_X}^{n}(X;k_x,\Omega_X^{n})\otimes L'
  \longrightarrow
  \operatorname{Ext}_{\Osh_X}^{n}(X;\Osh_X,\Omega_X^{n})\otimes L'
\end{equation}
deduzido do homomorfismo sobre os $\operatorname{Ext}^{n}$ associado a
$\Osh_X\longrightarrow k_x$:
\begin{equation}\tag{5.5}
  \operatorname{Ext}_{\Osh_X}^{n}(X;k_x,\Omega_X^{n})
  \longrightarrow
  \operatorname{Ext}_{\Osh_X}^{n}(X;\Osh_X,\Omega_X^{n})
  =H^{n}(X,\Omega_X^{n}).
\end{equation}

\src{183}
Como (5.3 bis) é um isomorfismo,\footnote{O texto impresso dá uma
segunda fórmula « (5.3) » na página 182 e remete aqui a ela; o identificador
« (5.3 bis) » distingue esta fórmula da primeira (5.3).} o mesmo ocorre
com (5.4), e portanto também com (5.5). Como $s_{(x)}$ é uma base de
$\operatorname{Ext}_{\Osh_X}^{n}(X;k_x,\Omega_X^{n})$ em virtude de (4.2), se
segue-se que sua imagem $\varepsilon_X$ é uma base de
$H^{n}(X,\Omega_X^{n})$.

Fica por demonstrar o enunciado (D), que resultará de maneira puramente formal dos fatos
elementares resumidos nos lemas seguintes. Neles se supõe que $X$ é uma
parte fechada (singular ou não) do espaço projetivo $P$ de dimensão $r$.
Utiliza-se a notação $\Osh_P(m)$ para o feixe sobre $P$ denotado por $\Osh(m)$ em
[3], e a notação $\Osh_X(m)$ para o feixe análogo sobre $X$.

\result{LEMA}{2} O enunciado (D) é verdadeiro se $X=P$ e $F=\Osh_P(m)$.

Este lema verifica-se mediante um cálculo direto. O cálculo explícito dos
$H^{i}(P,\Osh_P(m))$ encontra-se em [3], mas pode ser feito de maneira mais
elementar. O cálculo do produto cup
\[
  H^{i}(P,\Osh_P(m))\times H^{j}(P,\Osh_P(m'))
  \longrightarrow H^{i+j}(P,\Osh_P(m+m'))
\]
necessário para explicitar o acoplamento (5.1) não apresenta
dificuldades.\footnote{O texto impresso diz $\Osh_P(r+r')$; os dois fatores
$\Osh_P(m)$ e $\Osh_P(m')$ exigem $\Osh_P(m+m')$.}

\result{LEMA}{3} Todo feixe algébrico coerente $F$ sobre $X$ é isomorfo a
um feixe quociente de um feixe $\Osh_X(-m)^{k}$, onde se pode supor que $m$ é
tão grande quanto se queira.

Resulta do fato de que $F\otimes\Osh_X(m)$ está « gerado por suas seções » para
$m$ grande, cf. [3].

\result{LEMA}{4} Seja $i>0$. Então
$H^{r-i}(P,\Osh_P(-m))=0$ para $m$ grande; e para todo feixe algébrico
coerente $B$ sobre $X$, tem-se
\[
  \operatorname{Ext}_{\Osh_X}^{i}(X;\Osh_X(-m),B)=0
\]
para $m$ grande.

O primeiro enunciado resulta dos cálculos explícitos mencionados anteriormente; para o
segundo, observa-se que se tem um isomorfismo
\[
  \operatorname{Ext}_{\Osh_X}^{i}(X;\Osh_X(-m),B)
  =H^{i}(X,B\otimes\Osh_X(m))
\]
(proposição 1, corolário 1), de onde se segue a conclusão em virtude de um resultado bem
conhecido de [3]. Combinando os dois lemas anteriores, obtém-se o

\result{COROLÁRIO}{} Seja $i>0$. Então o funtor
$F\longmapsto H^{r-i}(P,F)$ sobre a categoria dos feixes algébricos coerentes
sobre $P$ é apagável, e o mesmo ocorre com o funtor
$\operatorname{Ext}_{\Osh_X}^{i}(X;F,B)$ sobre a categoria dos feixes
algébricos coerentes sobre $X$.

\src{184}
\result{LEMA}{5} Sejam $A$ e $B$ dois feixes algébricos coerentes sobre
$X$, e ponhamos $A(m)=A\otimes\Osh_X(m)$. Então, para $m$ suficientemente grande,
o homomorfismo canônico
\[
  \operatorname{Ext}_{\Osh_X}^{1}(X;A(-m),B)
  \longrightarrow H^{0}\bigl(X,\Extsh_{\Osh_X}^{1}(A(-m),B)\bigr)
  =H^{0}\bigl(X,\Extsh_{\Osh_X}^{1}(A,B)(m)\bigr)
\]
é um isomorfismo.

Isto resulta imediatamente da sequência espectral da proposição 1 aplicada a
$A(-m)$ e $B$, pois se terá
\[
  E_{2}^{p,q}(A(-m),B)
  =H^{p}\bigl(X,\Extsh_{\Osh_X}^{q}(A(-m),B)\bigr)
  =H^{p}\bigl(X,\Extsh_{\Osh_X}^{q}(A,B)(m)\bigr),
\]
que é nulo para $p>0$ e $m$ grande.

Demonstremos agora (D) no caso em que $X=P$. Provemos primeiro que (5.2 bis)
é um isomorfismo para $p=n$; como os dois membros são então
funtores exatos pela esquerda (pois $H^{r+1}(P,F)=0$), do lema 3 resulta
que basta demonstrar a afirmação para $F=\Osh_P(-m)$, mas esta se encontra
contida no lema 2. Como os homomorfismos (5.2 bis) são funtoriais e
compatíveis com os operadores de conexão, e como para $p<n$ os dois membros de
(5.2 bis) são funtores apagáveis em $F$ (corolário do lema 4), segue-se
então, mediante um raciocínio padrão, que (5.2 bis) é um isomorfismo
para todo $p$. Isto demonstra o teorema de dualidade para o espaço projetivo.
Suponhamos agora que $X$ é qualquer, mas não singular. Pelo teorema de
dualidade para $P$, tem-se um isomorfismo
\[
  H^{n}(X,F)'=H^{n}(P,F)'
  =\operatorname{Ext}_{\Osh_P}^{r-n}(P;F,\Omega_P^{r}).
\]
Em virtude do lema 1 (n\textsuperscript{o} 4), o último membro identifica-se com
\[
  \operatorname{Hom}_{\Osh_P}(P;F,\omega_X^{n})
  =\operatorname{Hom}_{\Osh_X}(X;F,\Omega_X^{n})
  =\operatorname{Ext}_{\Osh_X}^{0}(X;F,\Omega_X^{n}),
\]
de onde se obtém um isomorfismo
\begin{equation}\tag{5.6}
  H^{n}(X,F)'=\operatorname{Hom}_{\Osh_X}(X;F,\Omega_X^{n})
  =\operatorname{Ext}_{\Osh_X}^{0}(X;F,\Omega_X^{n}).
\end{equation}
Tomando $F=\Omega_X^{n}$, obtém-se um isomorfismo
\begin{equation}\tag{5.7}
  \gamma:H^{n}(X,\Omega_X^{n})\xrightarrow{\sim}k.
\end{equation}

\src{185}
Verifica-se que o isomorfismo (5.6) não é outro senão (5.2 bis) para $p=n$, quando
se toma $L=k$ graças a (5.7). Por conseguinte, (5.2 bis) é um isomorfismo para
$p=n$. Para demonstrar que é um isomorfismo para todo $p$, basta novamente
demonstrar que, para $p<n$, os dois membros de (5.2 bis) são funtores
apagáveis em $F$, e a fortiori (tendo em conta o lema 3) que ambos os membros
são nulos quando se toma $F=\Osh_X(-m)$ com $m$ grande. Ora, para o primeiro
membro, isto é verdadeiro em virtude do lema 4, e para o segundo membro escreve-se,
utilizando o teorema de dualidade sobre $P$:
\[
  H^{p}(X,\Osh_X(-m))'
  =\operatorname{Ext}_{\Osh_P}^{r-p}
    (P;\Osh_X(-m),\Omega_P^{r}).
\]
O segundo membro é nulo para $p<n$ e $m$ grande, como resulta do lema 5
(onde se toma $X=P$) e do fato de que $\Osh_X$, como feixe algébrico
coerente sobre $P$, tem dimensão cohomológica $\leq r-n$ (pois $X$ é não
singular), de onde
\[
  \Extsh_{\Osh_P}^{r-p}(\Osh_X,\Omega_P^{r})=0
  \qquad\text{para }p<n.
\]

\fsection{6}{O teorema de dualidade para as variedades singulares}

Seja $X$ uma parte fechada de dimensão $n$ do espaço projetivo $P$ de
dimensão $r$. A fórmula (5.6) deve ser escrita aqui
\begin{equation}\tag{6.1}
  H^{n}(X,F)'
  \simeq\operatorname{Hom}_{\Osh_X}(X;F,\omega_X^{n})
  =\operatorname{Ext}_{\Osh_X}^{0}(X;F,\omega_X^{n}),
\end{equation}
onde se pôs
\begin{equation}\tag{6.2}
  \omega_X^{n}=E^{0}(\Osh_X)
  =\Extsh_{\Osh_P}^{r-n}(\Osh_X,\Omega_P^{r}).
\end{equation}
Como foi dito na proposição 5, o feixe assim introduzido não depende,
na realidade, da imersão escolhida de $X$ em uma variedade não singular $P$.
Tomando, em (6.1), $F=\omega_X^{n}$, obtém-se
\begin{equation}\tag{6.2 bis}
  H^{n}(X,\omega_X^{n})'
  \simeq\operatorname{Hom}_{\Osh_X}(X;\omega_X^{n},\omega_X^{n}),
\end{equation}
\footnote{O texto impresso numera também esta fórmula « (6.2) »; o sufixo
« bis » distingue as duas ocorrências.} de onde resulta a existência de um elemento
distinguido em $H^{n}(X,\omega_X^{n})'$, correspondente ao homomorfismo
identidade de $\omega_X^{n}$ em si mesmo:
\begin{equation}\tag{6.3}
  \gamma:H^{n}(X,\omega_X^{n})\longrightarrow k.
\end{equation}

\src{186}
Consideremos então os acoplamentos definidos pela composição dos $\operatorname{Ext}$:
\begin{equation}\tag{6.4}
  H^{p}(X,F)\times
  \operatorname{Ext}_{\Osh_X}^{n-p}(X;F,\omega_X^{n})
  \longrightarrow H^{n}(X,\omega_X^{n}),
\end{equation}
e componhamo-los com o homomorfismo $\gamma$ de (6.3); obtêm-se
homomorfismos funtoriais, compatíveis com os operadores de bordo:
\begin{equation}\tag{6.5}
  \operatorname{Ext}_{\Osh_X}^{n-p}(X;F,\omega_X^{n})
  \longrightarrow H^{p}(X,F)'
\end{equation}
(que generalizam (5.2)). Verifica-se que para $p=n$ se obtém assim
o isomorfismo (6.1). Posto isto, tem-se o

\result{TEOREMA}{3 bis} As quatro condições seguintes sobre $X$ são
equivalentes para um inteiro dado $k\geq0$:
\begin{enumerate}
  \item[\textit{i.}] O homomorfismo funtorial (6.5) é um isomorfismo para
    $n-k\leq p\leq n$.
  \item[\textit{ii.}] Tem-se $H^{p}(X,\Osh_X(-m))=0$ para $m$ grande e para
    $n-k\leq p<n$.
  \item[\textit{iii.}] O funtor $H^{p}(X,F)$ sobre a categoria dos feixes
    algébricos coerentes sobre $X$ é apagável para $n-k\leq p<n$.
  \item[\textit{iv.}] Tem-se
    $E^{i}(\Osh_X)=\Extsh_{\Osh_P}^{r-n+i}(\Osh_X,\Omega_P^{r})=0$
    para $0<i\leq k$.
\end{enumerate}

\noindent\textbf{DEMONSTRAÇÃO.} --- \textit{i.}$\Rightarrow$\textit{ii.} Em
virtude do lema 4, \textit{ii.}$\Rightarrow$\textit{iii.} em virtude do lema 3,
\textit{iii.}$\Rightarrow$\textit{i.} mediante um raciocínio padrão bem conhecido,
tendo em conta que os dois membros de (6.5) são então funtores apagáveis
para $n-k\leq p<n$ (o primeiro o é em virtude do lema 4). Por fim, demonstremos
\textit{ii.}$\Leftrightarrow$\textit{iv.} Isto resulta do corolário da
proposição seguinte:

\result{PROPOSIÇÃO}{6} Seja $F$ um feixe algébrico coerente sobre $X$, e
seja $i$ um inteiro. Então, para $m$ suficientemente grande, tem-se um isomorfismo:
\begin{equation}\tag{6.6}
  H^{i}(X,F(-m))'
  \simeq H^{0}\bigl(X,E^{n-i}(F)(m)\bigr),
\end{equation}
onde se põe (compare-se n\textsuperscript{o} 3, proposição 5):
\begin{equation}\tag{6.7}
  E^{j}(F)=\Extsh_{\Osh_P}^{r-n+j}(F,\Omega_P^{r}).
\end{equation}
Com efeito, em virtude do teorema de dualidade sobre $P$, o primeiro membro de (6.6)
é isomorfo a
$\operatorname{Ext}_{\Osh_P}^{r-i}(P;F(-m),\Omega_P^{r})$, pelo que (6.6) resulta
do lema 5.

\src{187}
\result{COROLÁRIO}{} Para que se tenha $H^{i}(X,F(-m))=0$ para $m$ grande, é
necessário e suficiente que $E^{n-i}(F)=0$.

Lembremos que os $E^{j}(F)$ não dependem, na realidade, da imersão projetiva
considerada. A condição do corolário é de natureza puramente local; por conseguinte,
se ela se verifica para $F$, verifica-se para todo feixe localmente
isomorfo a um feixe $F^{n}$. Em particular, se esta condição se verifica
para $\Osh_X$, verifica-se para todo feixe algébrico coerente localmente
livre. Este é, por exemplo, o caso para todo $i<n$ se $X$ é não singular, para
$i=0$ se nenhuma componente de $X$ se reduz a um só ponto, e para $i=0,1$
se $X$ é normal e todas as suas componentes têm dimensão $>1$ (cf.
[3]).\footnote{O texto impresso diz aqui « $S$ é normal »; o parágrafo
refere-se apenas à variedade $X$.} Para que isto ocorra para todo
$i<k$, é necessário e suficiente, por definição, que os anéis locais
$\Osh_x$ ($x\in X$) tenham « codimensão homológica » $\geq k$ (cf. [4]
para esta noção). Se $k=n$, isto significa, em virtude do teorema 3 bis, que o
teorema de dualidade é verdadeiro para $X$, isto é, que (6.5) é um isomorfismo para
todo $p$ e todo $F$. Podem-se dar numerosas condições equivalentes sobre
os anéis locais $\Osh_x$ para que isto ocorra (Nagata), por exemplo, que
satisfaçam o teorema de equidimensionalidade de Cohen--Macaulay. Assim ocorre,
por exemplo, se localmente $X$ é uma « interseção completa » em uma
variedade ambiente não singular.

\fsection{7}{A dualidade de Poincaré}

Seja $X$ uma variedade projetiva não singular de dimensão $n$. Então
\[
  H^{*}(X)=H^{*}(X,\Omega_X^{*})
\]
é uma álgebra anticonmutativa bigraduada de dimensão finita, que graduamos
pelo grau total, de modo que $H^{p,q}(X)=H^{p}(X,\Omega_X^{q})$ tem grau
$p+q$. Os graus de $H^{*}(X)$ estão compreendidos entre $0$ e $2n$. Em virtude do
teorema 2 e do corolário do teorema 3, $H^{*}(X)$ é uma « álgebra de
Poincaré » de dimensão $2n$, isto é, $H^{2n}(X)$ está provido de um isomorfismo com
o corpo base $k$, e o produto
\[
  H^{m}(X)\times H^{2n-m}(X)\longrightarrow H^{2n}(X)=k
\]
é uma dualidade entre $H^{m}(X)$ e $H^{2n-m}(X)$. Além disso, se $Y$ é outra
variedade projetiva não singular, a fórmula de Künneth para os
feixes algébricos coerentes dá
\begin{equation}\tag{7.1}
  H^{*}(X\times Y)=H^{*}(X)\otimes H^{*}(Y),
\end{equation}
isomorfismo compatível com as estruturas de álgebras de Poincaré. Além disso,
$H^{*}(X)$ é, como álgebra comutativa, um funtor contravariante em
$X$, pois um morfismo $f:Y\longrightarrow X$ define de maneira evidente um
homomorfismo de álgebras graduadas

\src{188}
\begin{equation}\tag{7.2}
  f^{*}:H^{*}(X)\longrightarrow H^{*}(Y).
\end{equation}
Como se trata de álgebras de Poincaré, obtém-se por transposição um
homomorfismo de espaços vetoriais:
\begin{equation}\tag{7.3}
  f_{*}:H^{*}(Y)\longrightarrow H^{*}(X).
\end{equation}
Viu-se no n\textsuperscript{o} 4 que o efeito de $f^{*}$ sobre as classes de
cohomologia correspondentes a ciclos não singulares se interpreta
geometricamente tomando as classes de cohomologia que correspondem a suas
imagens inversas. No caso atual, importa demonstrar que (7.3) corresponde
analogamente à operação « imagem direta » de ciclos. Isto resulta, ao menos sob
condições adequadas de não singularidade, do caso particular seguinte:

\result{TEOREMA}{4} Se $f$ é a inclusão de uma subvariedade não
singular $Y^{m}$ de $X^{n}$ em $X^{n}$, então, denotando por $1_Y$
o elemento unidade de $H(Y)$, tem-se
\begin{equation}\tag{7.4}
  f_{*}(1_Y)=P_X(Y),
\end{equation}
onde o segundo membro é a classe de cohomologia em $X$ associada a $Y$.

Esta fórmula equivale a
\begin{equation}\tag{7.4 bis}
  \bigl\langle\xi^{m,m},P_X(Y)\bigr\rangle\varepsilon_Y
  =f^{*}(\xi^{m,m})
  \qquad
  \bigl(\xi^{m,m}\in H^{m}(X,\Omega_X^{m})\bigr),
\end{equation}
\footnote{O texto impresso diz $f_{*}(\xi^{m,m})$. Visto que
$\xi^{m,m}\in H^{m}(X,\Omega_X^{m})$ e que ambos os membros pertencem a
$H^{m}(Y,\Omega_Y^{m})$, a aplicação requerida é $f^{*}$.} onde
$\varepsilon_Y$ é o elemento fundamental em $H^{m}(Y,\Omega_Y^{m})$, e
constitui, no caso das variedades projetivas não singulares, uma nova
definição da classe de cohomologia associada a $Y$. Para demonstrar o
teorema 4, interpreta-se, graças ao teorema 3, o transposto do homomorfismo
\[
  H^{m}(X,\Omega_X^{m})
  \longrightarrow H^{m}(Y,\Omega_Y^{m})
  =H^{m}(X,\Omega_Y^{m})
\]
como o homomorfismo
\begin{equation}\tag{7.5}
\begin{split}
  H^{n-m}(X,\Omega_X^{n-m})
  &\simeq
  \operatorname{Ext}_{\Osh_X}^{n-m}(X;\Omega_X^{m},\Omega_X^{n})\\
  &\longleftarrow
  \operatorname{Ext}_{\Osh_X}^{n-m}(X;\Omega_Y^{m},\Omega_X^{n})
  \simeq
  \operatorname{Hom}_{\Osh_X}(X;\Omega_Y^{m},\Omega_Y^{m}).
\end{split}
\end{equation}
Verifica-se que o elemento $1_Y$ do dual de $H^{m}(Y,\Omega_Y^{m})$ identifica-se
com o elemento do membro da direita correspondente ao endomorfismo identidade de
$\Omega_Y^{m}$, e por outro lado que a imagem deste elemento em
$H^{n-m}(X,\Omega_X^{n-m})$ é precisamente $P_X(Y)$. Estas compatibilidades, que poderiam
ter sido dadas já no n\textsuperscript{o} 4, podem ser enunciadas, e são
na realidade válidas, para variedades não singulares quaisquer; a segunda

\src{189}
resulta, por exemplo, da comutatividade do diagrama seguinte de endomorfismos
canônicos:
\begin{equation}\tag{7.6}
\resizebox{\textwidth}{!}{$
\left\{
\begin{array}{ccccc}
\operatorname{Ext}_{\Osh_X}^{n-m}(X;\Omega_X^{m},\Omega_X^{n})
&\longleftarrow&
\operatorname{Ext}_{\Osh_X}^{n-m}(X;\Omega_Y^{m},\Omega_X^{n})
&\simeq&
\operatorname{Hom}_{\Osh_X}(X;\Omega_Y^{m},\Omega_Y^{m})\\[3pt]
\Big\downarrow&&&&\Big\downarrow\\[3pt]
\operatorname{Ext}_{\Osh_X}^{n-m}(X;\Osh_X,\Omega_X^{n-m})
&\longleftarrow&
\operatorname{Ext}_{\Osh_X}^{n-m}(X;\Osh_Y,\Omega_X^{n-m})
&\simeq&
\operatorname{Hom}_{\Osh_X}(X;\Omega_Y^{m},\Omega_Y^{m})
\end{array}
\right.
$}
\end{equation}

Obtém-se assim um equivalente exato do formalismo da dualidade de Poincaré
sobre as variedades orientadas compactas. Em particular, o teorema 4 permite
determinar a classe de cohomologia associada à diagonal de $X\times X$. Mediante
um raciocínio bem conhecido, deduz-se, por exemplo, uma fórmula de Lefschetz:

\result{TEOREMA}{5} Seja $f$ um endomorfismo de uma variedade projetiva não
singular $X$, tal que os pontos fixos de $f$ tenham multiplicidade 1.
Então o número desses pontos fixos é congruente, módulo a característica de
$k$, com a soma alternada dos traços dos endomorfismos dos $H^{i}(X)$ definidos
por $f$.

A restrição sobre $f$ que foi necessário impor deve-se às dificuldades mencionadas na
observação do n\textsuperscript{o} 4. Não obstante, observar-se-á que a fórmula de
Lefschetz continua verdadeira se $f$ é « homotópico » a um endomorfismo cujos
pontos fixos têm todos multiplicidade 1.

\fsection{8}{Generalização do teorema de dualidade}

Seja $X$ uma variedade algébrica não singular tal que todo feixe
algébrico coerente $F$ sobre $X$ seja isomorfo a um quociente de um feixe
algébrico coerente localmente livre (o que ocorre se $X$ é localmente
fechado em um espaço projetivo). Então todo feixe algébrico coerente $F$ sobre
$X$ admite uma resolução finita $L$ mediante feixes localmente livres, e para
duas resoluções desta classe sempre se pode encontrar uma terceira e um
homomorfismo desta nas duas primeiras compatível com as
aumentações. Analogamente, se $L$ é uma resolução finita localmente livre de
$F$, e se é dado um homomorfismo $F'\longrightarrow F$, então existe
uma resolução finita localmente livre $L'$ de $F'$ e um homomorfismo
$L'\longrightarrow L$ compatível com $F'\longrightarrow F$, que se pode
supor sobrejetivo se $F'\longrightarrow F$ o é. Isto permite definir, dados
dois inteiros $r,s\geq0$, dois multifuntores cohomológicos nos
argumentos $A_1,\ldots,A_r$; $B_1,\ldots,B_s$,

\src{190}
que variam na categoria dos feixes algébricos coerentes sobre $X$, e cujos
valores estão na categoria dos feixes algébricos coerentes sobre $X$,
resp. na categoria dos módulos sobre $H^{0}(X,\Osh_X)$, mediante as fórmulas:
\begin{equation}\tag{8.1}
\resizebox{\textwidth}{!}{$
\left\{
\begin{aligned}
\underline{T}_{r}^{s*}(A_1,\ldots,A_r;B_1,\ldots,B_s)
  &={}H^{*}\!\left(
    \Homsh_{\Osh_X}\!\left(
      L(A_1)\otimes\cdots\otimes L(A_r),
      L(B_1)\otimes\cdots\otimes L(B_s)
    \right)\right),\\[3pt]
T_{r}^{s*}(A_1,\ldots,A_r;B_1,\ldots,B_s)
  &={}\mathbf R^{*}\Gamma\!\left(
    \Homsh_{\Osh_X}\!\left(
      L(A_1)\otimes\cdots\otimes L(A_r),
      L(B_1)\otimes\cdots\otimes L(B_s)
    \right)\right).
\end{aligned}
\right.
$}
\end{equation}
Nesta fórmula, $L(F)$ designa uma resolução finita localmente livre do
feixe algébrico coerente $F$, e $\mathbf R^{*}\Gamma(K)$ designa
a hipercohomologia do espaço $X$ com respeito ao complexo de feixes $K$. Se
$r$ ou $s$ é nulo, substitui-se o produto tensorial dos $L(A_i)$, resp. dos
$L(B_j)$, por $\Osh_X$. Em particular, $\underline{T}_{0}^{0}$ e $T_{0}^{0}$
são funtores graduados de 0 argumentos, $\underline{T}_{0}^{0}$ se reduz ao
grau 0, onde é o feixe $\Osh_X$, enquanto que $T_{0}^{0}$ é idêntico a
$H^{*}(X,\Osh_X)$. O fato de que os segundos membros de (8.1) não dependam
das resoluções escolhidas é, em qualquer caso, evidente para a primeira linha
(pois a questão é então local), e para a segunda resulta das observações
gerais precedentes, tendo em conta a sequência espectral de
hipercohomologia do complexo de feixes
\[
  K=\Homsh_{\Osh_X}\bigl(L(A_1)\otimes\cdots,
      L(B_1)\otimes\cdots\bigr),
\]
que converge à hipercohomologia de $X$ com respeito a $K$ e cujo termo
inicial é $H^{p}(X,H^{q}(K))$, isto é,
\begin{equation}\tag{8.2}
  E_{2}^{p,q}
  =H^{p}\!\left(X,
    (\underline{T}_{r}^{s})^{(q)}
      (A_1,\ldots,A_r;B_1,\ldots,B_s)\right).
\end{equation}
Vê-se então que esta sequência espectral tampouco depende das resoluções
escolhidas; seu limite é o primeiro membro de (8.1). Definem-se
facilmente os operadores de conexão relativos aos diversos argumentos $A_i$, $B_j$,
observando que toda sequência exata
$0\longrightarrow F'\longrightarrow F\longrightarrow F''\longrightarrow0$
pode ser resolvida mediante uma sequência exata de complexos finitos localmente livres.

No sistema dos funtores $\underline{T}_{r}^{s*}$, resp.
$T_{r}^{s*}$, definem-se operações análogas às do cálculo tensorial, cuja
definição é imediata a partir das fórmulas definidoras (8.1). Assim, tem-se
uma composição (que generaliza a considerada no n\textsuperscript{o} 2):
\begin{equation}\tag{8.3}
\resizebox{\textwidth}{!}{$
\begin{aligned}
&T_{r}^{s*}(A_1,\ldots,A_r;B_1,\ldots,B_s)
\times T_{r'}^{s'*}(A'_1,\ldots,A'_{r'};B'_1,\ldots,B'_{s'})\\
&\qquad\longrightarrow
T_{r+r'}^{s+s'*}
(A_1,\ldots,A_r,A'_1,\ldots,A'_{r'};
 B_1,\ldots,B_s,B'_1,\ldots,B'_{s'}),
\end{aligned}
$}
\end{equation}
que satisfaz as propriedades evidentes de associatividade e de compatibilidade com
os homomorfismos funtoriais, os homomorfismos de conexão e as sequências
espectrais. Do mesmo modo, têm-se operações de simetria cuja explicitação se deixa ao
leitor. Tem-se, além disso, uma operação de contração cada vez que
um dos argumentos $A_i$

\src{191}
é igual a um dos argumentos $B_j$:
\begin{equation}\tag{8.4}
\resizebox{\textwidth}{!}{$
\begin{aligned}
&T_{r}^{s*}(A_1,\ldots,A_{i-1},C,A_{i+1},\ldots,A_r;
 B_1,\ldots,B_{j-1},C,B_{j+1},\ldots,B_s)\\
&\qquad\longrightarrow
T_{r-1}^{s-1*}(A_1,\ldots,\widehat{A_i},\ldots,A_r;
 B_1,\ldots,\widehat{B_j},\ldots,B_s).
\end{aligned}
$}
\end{equation}
Além disso, se um argumento $A_i$ é um feixe localmente livre, pode
ser suprimido, com a condição de substituir um dos $A_j$ ($j\neq i$) por
$A_j\otimes A_i$, ou um dos $B_k$ por $B_k\otimes A_i'$ (onde
$A_i'=\Homsh_{\Osh_X}(A_i,\Osh_X)$), e tem-se uma regra de cálculo análoga no
caso em que um dos argumentos $B_i$ é localmente livre. Em particular, sempre
se pode suprimir um argumento que seja idêntico a $\Osh_X$. Se todos os
argumentos são localmente livres, salvo no máximo um dos argumentos $B_i$, a regra
que acaba de ser enunciada dá um isomorfismo funtorial
\begin{equation}\tag{8.5}
  T_{r}^{s*}(A_1,\ldots,A_r;B_1,\ldots,B_s)
  =H^{*}\bigl(X,A_1'\otimes\cdots\otimes A_r'
    \otimes B_1\otimes\cdots\otimes B_s\bigr)
\end{equation}
(pois se reduz ao caso em que $r=0$, $s=1$, onde isto é imediato; também se
pode utilizar diretamente a sequência espectral de termo inicial (8.2)). Para os
$\underline{T}_{r}^{s}$ se definem operações correspondentes às anteriores.
As relações entre as diversas operações assim
introduzidas são as mesmas que para as operações análogas do cálculo
tensorial.

Seja $n$ a dimensão de $X$. Aplicando sucessivamente uma composição
tensorial (8.3) e contrações (8.4) sobre os argumentos repetidos, obtém-se
um acoplamento
\begin{equation}\tag{8.6}
\resizebox{\textwidth}{!}{$
  (T_{r}^{s})^{p}(A_1,\ldots;B_1,\ldots)
  \times
  (T_{s}^{r})^{n-p}(B_1,\ldots;A_1,\ldots,A_r\otimes\Omega_X^{n})
  \longrightarrow H^{n}(X,\Omega_X^{n}).
$}
\end{equation}

\result{TEOREMA}{6} Se $X$ é uma variedade projetiva não singular, então
os acoplamentos (8.6) são dualidades.

Isto resulta de maneira puramente formal do corolário do teorema 3. Com efeito,
deste corolário resulta facilmente que, se $K$ é um complexo de feixes
algébricos coerentes localmente livres, então a hipercohomologia de $X$ com
respeito a $K$ está em dualidade com a hipercohomologia de $X$ com respeito a
$K'\otimes\Omega_X^{n}$ mediante os acoplamentos naturais
\begin{equation}\tag{8.7}
  \mathbf R^{p}\Gamma(K)
  \times\mathbf R^{n-p}\Gamma(K'\otimes\Omega_X^{n})
  \longrightarrow\mathbf R^{n}\Gamma(\Omega_X^{n})
  =H^{n}(X,\Omega_X^{n}).
\end{equation}
Isto se vê utilizando a sequência espectral de termo inicial
$H^{p}(X,H^{q}(K))$\footnote{O texto impresso diz
$H^{p}(H^{q}(X,K))$; o termo inicial é o da sequência espectral (8.2),
isto é, $H^{p}(X,H^{q}(K))$.} e a sequência espectral análoga para
$K'\otimes\Omega_X^{n}$. Do resultado anterior se deduz o teorema 6
utilizando a definição (8.1).

\src{192}
\subsection*{OBSERVAÇÕES}

\noindent 1\textsuperscript{o} Para as definições anteriores ao teorema 6,
não era necessário que $X$ fosse não singular, pois não era
indispensável trabalhar apenas com resoluções finitas. Mas se $X$
é singular, já não se poderá afirmar, a priori, que os
$(\underline{T}_{r}^{s})^{p}(A_1,\ldots;B_1,\ldots)$ sejam feixes
coerentes, pois no complexo de feixes
\[
  \Homsh_{\Osh_X}\bigl(L(A_1)\otimes\cdots,L(B_1)\otimes\cdots\bigr)
\]
haverá uma infinidade de componentes de um grau total dado.

\medskip
\noindent 2\textsuperscript{o} Verifica-se facilmente que, nas fórmulas
(8.1), pode-se substituir um dos $L(B_i)$ por $B_i$. Tendo em conta a
proposição 3, isto demonstra que se tem
\begin{equation}\tag{8.8}
\left\{
\begin{aligned}
  \underline{T}_{1}^{1*}(A;B)&=\Extsh_{\Osh_X}^{*}(A,B),\\
  T_{1}^{1*}(A;B)&=\operatorname{Ext}_{\Osh_X}^{*}(X;A,B).
\end{aligned}
\right.
\end{equation}
Em particular, tomando em (8.6) $r=s=1$ e $A_1=\Osh_X$, recupera-se o
teorema 3. A fórmula (8.8) implica também
\[
  T_{0}^{1*}(B)=H^{*}(X,B),
  \qquad
  T_{1}^{0*}(A)=\operatorname{Ext}_{\Osh_X}^{*}(X;A,\Osh_X).
\]

\medskip
\noindent 3\textsuperscript{o} Vê-se na fórmula (8.1) que, em geral, os
funtores $\underline{T}_{r}^{s*}$ e $T_{r}^{s*}$ têm componentes de
graus positivos e negativos. Utilizando a observação 2, vê-se que, se a
dimensão de $X$ é $n$, então as componentes não nulas de
$\underline{T}_{r}^{s*}$ estão compreendidas entre $-(s-1)n$ e $rn$, se $s>0$,
entre $0$ e $rn$, se $s=0$, e as componentes não nulas de $T_{r}^{s*}$
estão compreendidas entre $-(s-1)n$ e $(r+1)n$, se $s>0$, entre $0$ e $(r+1)n$,
se $s=0$ (e inclusive, salvo erro, se $r>0$, entre $-(s-1)n$ e $rn$, resp. entre
$0$ e $rn$).

\section*{BIBLIOGRAFIA}

\noindent[1] Cartier (Pierre). --- Les groupes
$\operatorname{Ext}^{s}(A,B)$, Séminaire A. Grothendieck : Algèbre homologique,
t. 1, 1957, n\textsuperscript{o} 3.

\noindent[2] Grothendieck (Alexandre). --- Sur quelques points d'algèbre
homologique, \emph{Tohoku Math. J.}, t. 9, 1957, p. 119--183.

\noindent[3] Serre (Jean-Pierre). --- Faisceaux algébriques cohérents,
\emph{Annals of Math.}, t. 61, 1955, p. 197--278.

\noindent[4] Serre (Jean-Pierre). --- Sur la dimension homologique des anneaux
et des modules noethériens, \emph{Proc. Intern. Symp. on Alg. Number Theory}
[1955, Tokyo et Nikko]. --- Tokyo, Science Council of Japan, 1956 ;
p. 175--189.

\clearpage
\src{193}
\section*{ADENDA}

As dificuldades assinaladas na observação da página 13 estão atualmente
completamente resolvidas, graças a uma extensão dos teoremas de dualidade a
variedades quaisquer (com singularidades arbitrárias). Para esta teoria, que pode
ser formulada no quadro geral da « teoria dos esquemas », remetemos
aos « Éléments de Géométrie algébrique » de J. Dieudonné e A. Grothendieck,
atualmente em preparação. Encontrar-se-ão algumas indicações em:

\begin{quote}
Grothendieck (A.). --- Cohomology theory of algebraic varieties,
\emph{Congrès int. des Mathématiciens}, 1958, Edinburgh (de próxima publicação).
\end{quote}

\begin{flushright}
[Abril de 1959]
\end{flushright}

\begin{center}\rule{36mm}{0.5pt}\end{center}
